\section{Geometric preparation with the physical matter margin}
\label{ch:prep:section}

Fix a strict rest datum from Proposition~\ref{ch:rest:reduction}, denoted
by $(N,g,K,\alpha_1,\alpha_2)$, with obstacle $S_0$, energy $E_*>0$,
and $A_*=a_g(S_0)$. Retain its geometry and end estimates; in
particular, $K$ has compact support and $\theta_+(S_0)<0$.
Its physical and neutral margins satisfy
\begin{equation}\label{ch:prep:margin}
 M_0=8\pi(\mu_m-|J_m|_g)>0,
 \qquad 8\pi(\mu-|J|_g)\geq M_0,
 \qquad M_0\geq c_*r_y^{-3-\delta_*}
 \quad\text{far out},
\end{equation}
where $c_*>0$ and $0<\delta_*<1$ are fixed for this datum.
Every choice and constant in this section concerns this fixed datum.

\subsection{The bounding regions and their orientations}

For a symmetric tensor $T$ write
$\Theta_T(\Sigma,\nu)=H_\Sigma(\nu)+\tr_\Sigma T$.
The geometric input is the following form of the three-dimensional
barrier and total trapped-region theorems
\cite[Section~2, Theorems~3.1 and~4.1, Definitions~7.1--7.2,
and Theorem~7.3]{AnderssonMetzger2009}.

\begin{theorem}[Barrier and total trapped-region theorem]
\label{ch:prep:region-theorem}
Let $(D,g,T)$ be smooth compact connected three-dimensional data with boundary
$\Gamma_-\sqcup\Gamma_+$.  Assume $\Gamma_+\ne\varnothing$ and
$\Theta_T(\Gamma_+)>0$ with its normal out of $D$.  The inner boundary
may be empty; when present it satisfies $\Theta_T(\Gamma_-)<0$ with
its normal into $D$.  All boundary parts may be disconnected.

Consider all smooth bounding domains containing a relative collar
of $\Gamma_-$, avoiding $\Gamma_+$, and having interior free boundary
with $\Theta_T\leq0$ for the normal out of the bounding domain.
If the union is nonempty, its closure is a maximal
closed bounding region whose interior frontier is smooth, compact,
embedded, stable, and satisfies $\Theta_T=0$.  It is itself the
closure of an admissible domain.  A nonempty strict inner barrier
ensures a nonempty union and an enclosing stable MOTS.
No energy condition or constraint equation is required for $T$.
\end{theorem}

Filling a complementary component incident on no designated outer
boundary deletes whole free-boundary components and preserves the
expansion of those retained.  This filling convention is understood
in Theorem~\ref{ch:prep:region-theorem} and throughout the construction.
Two identities fix the shifted tensors and all orientations:
\begin{equation}\label{ch:prep:shift-reversal}
 \Theta_{T-(c/2)g}=\Theta_T-c,
 \qquad \Theta_{-T}(\Sigma,-\nu)=-\Theta_T(\Sigma,\nu).
\end{equation}

Choose $R_0$ beyond $\supp K$ so that every coordinate sphere of
radius at least $R_0$ has positive outward mean curvature.  Every
compact bounding domain with $\Theta_{\pm K}\leq c\leq0$ is confined
inside this fixed radial range.  Indeed, at a free-boundary point
of maximum radius beyond $R_0$, the boundary lies inside the tangent
coordinate sphere with the same outward normal.  The local graph
comparison gives $H_{\partial D}\geq H_{S_r}>0$, whereas $K=0$.
This contradicts the expansion bound.  Consequently all families
below can be constructed on a single compact truncation, independently
of the sufficiently distant outer sphere.

For $\max_{S_0}\Theta_K(S_0)<c<0$, let $\mathcal B_c$ be the
\emph{full} closed maximal region for $\Theta_K\leq c$, with $S_0$
as required inner boundary.  Applying
Theorem~\ref{ch:prep:region-theorem} to $K-(c/2)g$ is legitimate:
$S_0$ is strictly negative and the distant sphere strictly positive
for the shifted expansion.  Thus $\mathcal B_c$ is nonempty, contains
a collar of $S_0$, and has smooth stable frontier at expansion $c$.
These regions increase with $c$.

Fix such a threshold $c_b$ and put $\mathcal B=\mathcal B_{c_b}$.
In its filled complement, designate both the full black frontier and a
distant sphere as \emph{outer} barriers.  There is no required inner
boundary.  For $c<0$, let $\mathcal W_c$ be the full closed maximal
region of compact bounding domains satisfying $\Theta_{-K}\leq c$
there.  On the reversed black frontier the expansion of
$-K-(c/2)g$ equals $-c_b-c>0$; it is positive on the distant
sphere as well.  Theorem~\ref{ch:prep:region-theorem} therefore applies
whenever the union is nonempty.  Its smooth stable frontier has
$\Theta_{-K}=c$ and positive separation from the black frontier.
If the union is empty, set $\mathcal W_c=\varnothing$.  The white
family increases with $c$, always with the black region held fixed.

We choose the two thresholds successively at points of Hausdorff
right continuity of these full families.  Here is the elementary
selection argument, including emptiness.  For any increasing closed
family $E_c$ in a compact metric space, choose a cap $L$ larger than
its diameter and put
$d_c(x)=\min\{L,\dist(x,E_c)\}$, taking $d_c\equiv L$ if
$E_c=\varnothing$.  At every point of a countable dense set, this
is a bounded nonincreasing function of $c$ and has at most countably
many right discontinuities.  Avoid their union.  The common
$1$-Lipschitz bound and a finite-net argument then give uniform
convergence $d_{c_j}\to d_c$ whenever $c_j\downarrow c$.  For
nonempty $E_c$ this gives Hausdorff convergence, since
$E_c\subset E_{c_j}$.  For empty $E_c$, uniform convergence to $L$
forces every sufficiently close $E_{c'}$, $c'>c$, to be empty:
otherwise its distance function has a zero.  Thus $c_b$, and then
$c_w$ for this fixed $\mathcal B$, can be chosen arbitrarily close
to zero with the stated continuity, including the empty-white case.

Put $\mathcal W=\mathcal W_{c_w}$ and let $X$ be the closure of the
component toward infinity of $N\setminus(\mathcal B\cup\mathcal W)$.
Its boundary is the entire adjacent frontier, a finite disjoint
union $B=B_b\sqcup B_w$ of complete black and white face components.
It is nonempty because the excluded side contains $S_0$ and its
collar.  Both full frontiers are compact and disjoint, so selection
creates no corners.  With $\nu$ pointing into $X$, set
$H=H_B(\nu)$ and $P_B=\tr_BK$.  Then
\begin{equation}\label{ch:prep:faces}
 H+P_B=c_b\quad\text{on }B_b,
 \qquad H-P_B=c_w\quad\text{on }B_w.
\end{equation}
Either face type may be absent from $B$; for example a white face
can separate a black face from infinity.  The full regions used to
choose the thresholds are retained in the preparation even when
some of their frontier components are not adjacent to $X$.

The restricted data on $X$ are smooth and complete up to $B$, with
the original single end.  Every enclosing cut $\Gamma$ for $B$ is
also a cut for $S_0$ after the discarded side is adjoined to its
inner region and bounded pockets are filled.  Therefore
\begin{equation}\label{ch:prep:enclosure}
 \Area_g(\Gamma)\geq A_*,
 \qquad
 \frac1{4\pi}\int_\Gamma\alpha_i
   =\frac1{4\pi}\int_B\alpha_i=Q_i\quad(i=1,2),
 \qquad (Q_1,Q_2)=(Q_E,Q_B).
\end{equation}
The flux equalities follow from closedness and Stokes' Theorem
between the cut and a distant sphere, all normals pointing toward
the end.  These assertions concern the whole frontier and do not
require any individual component to separate $S_0$ from infinity.

\subsection{Stable collars and a stronger offset bound}

\begin{lemma}[Outward leaf collars]\label{ch:prep:collar}
Each connected face has a smooth outward collar with leaf parameter
$s\geq0$, $s=0$ on the face, and $|\nabla s|$ bounded above and
below by positive constants.  Its leaves satisfy
$\Theta_K>c_b$ for black faces and $\Theta_{-K}>c_w$ for white
faces whenever $s>0$.  The collars may be chosen disjoint in $X$.
\end{lemma}

\begin{proof}
Let $\Sigma$ be a connected component of a full maximal frontier
at threshold $c$, for $T=K$ or $-K$, with its outward normal.
For normal graphs define
$\Phi(v)=\Theta_T(\Sigma(v))-c$ as a map
$C^{2,\alpha}(\Sigma)\to C^{0,\alpha}(\Sigma)$.
Its linearization $\mathcal L$ is the stability operator of
$T-(c/2)g$.  It has principal eigenvalue $\lambda\geq0$ and a
positive eigenfunction $\varphi$.  If $\lambda>0$, the graphs
$v=s\varphi$ have
$\Phi(s\varphi)=s\lambda\varphi+O(s^2)>0$ for small $s>0$.

If $\lambda=0$, the kernels of $\mathcal L$ and its adjoint are
one-dimensional, with positive generators $\varphi$ and
$\varphi^*$.  The bordered linear map
\begin{equation}\label{ch:prep:fredholm}
 (v,a)\longmapsto
 \left(\mathcal L v-a,\int_\Sigma v\,dA_g\right)
\end{equation}
is invertible.  In fact, solving $\mathcal L v-a=f$ first determines
$a=-\int_\Sigma\varphi^*f\,dA_g/\int_\Sigma\varphi^*\,dA_g$
by Fredholm compatibility, and adding a unique multiple of
$\varphi$ fixes the prescribed integral of $v$.  Elliptic estimates
give a bounded inverse.  The implicit-function theorem applied to
$(\Phi(v)-a,\int v)=(0,s)$ now yields constant-expansion leaves
with $a'(0)=0$ and $v'(0)=\varphi/\int_\Sigma\varphi\,dA_g>0$.
They foliate a short outward collar.  If $a(s)\leq0$ for any
positive small $s$, adjoining this collar and replacing just this
whole frontier component by its new leaf would strictly enlarge
the full maximal bounding region while preserving expansion
$\leq c$ on every free component.  The collar avoids all other
faces and designated barriers, so this is an admissible competitor,
contradicting maximality.  Thus $a(s)>0$.

Smooth elliptic regularity gives smooth leaves; positive graph
velocity gives the two bounds on $|\nabla s|$.  There are finitely
many faces.  Shortening the collars of those adjacent to $X$ makes
them disjoint and contained in $X$.
\end{proof}

\begin{proposition}[Prepared exterior with physical margin]
\label{ch:prep:prepared}
The thresholds and exterior above can be chosen with their stated
full-region right continuity and collars, and with the following
additional data.  There are constants $C_b\geq0$, $C_w\leq0$ and
a smooth compactly supported function $C$ satisfying
\begin{equation}\label{ch:prep:offsets}
 \begin{aligned}
 C_w\leq C\leq C_b,\qquad
 C&=C_b-k_Bs &&\text{near each black face},\\
 C&=C_w+k_Bs &&\text{near each white face},
 \end{aligned}
\end{equation}
where $k_B>0$ is constant on each individual collar.  Set
\begin{equation}\label{ch:prep:heights}
 b_-=c_b-C_b<0,\qquad b_+=-c_w-C_w>0.
\end{equation}
There is a smooth positive weight $\rho$, equal to a positive
constant times $r_y^{-4}$ on the end with differentiated symbol
bounds, such that, for every real $h\in[b_-,b_+]$ and every $x\in X$,
\begin{equation}\label{ch:prep:offset-smallness}
 |(h+C)\tr_gK|+|\nabla C|_g+\rho\leq\tfrac12M_0(x).
\end{equation}
The thresholds, full regions, exterior, collars, offset, assigned
heights, and weight are all fixed before any elliptic deformation
parameter.  After fixing the thresholds and collars, the offset can
have arbitrarily small $C^1$ norm.  The constants of an absent face
type remain valid bounds.  Equation~\eqref{ch:prep:enclosure} holds.
\end{proposition}

\begin{proof}
Take a fixed compact neighborhood $\mathcal K$ containing the
uniform confinement range and $\supp K$ in its interior, and put
$\eta_*=\min_{\mathcal K}M_0>0$ and
$T_*=\sup_{\mathcal K}|\tr_gK|$.  Restrict both negative thresholds
in advance to an interval so close to zero that
$\max\{|c_b|,|c_w|\}T_*<\eta_*/8$, while maintaining the strict
black inner-barrier inequality.  The countable-exception argument
permits the successive right-continuous choices there.  Construct
the regions and fix their disjoint collars inside $\mathcal K$.

On each collar choose a smooth cutoff $\zeta_B(s)$ in $[0,1]$,
equal to one near zero and zero near the far end.  Fix
$\kappa_B>0$ with $1-\kappa_Bs>0$ throughout the collar.
For a common amplitude $a_*>0$ define the black contribution to
$C$ as $a_*\zeta_B(s)(1-\kappa_Bs)$ and the white contribution
as its negative.  Extend by zero and sum on the disjoint collars.
Then $C_b=a_*$, $C_w=-a_*$, $k_B=a_*\kappa_B$, and
$\|C\|_{C^1}=O(a_*)$ with constants depending only on the fixed
collars and profiles.  For every assigned height,
\[
 |h+C|\leq\max\{|c_b|,|c_w|\}+2a_*.
\]
Choose $a_*$ small enough that
$2a_*T_*+\|\nabla C\|_\infty<\eta_*/8$.
The first two terms of Equation~\eqref{ch:prep:offset-smallness} are
then at most $\eta_*/4\leq M_0/4$ on $\mathcal K$, and vanish
outside it.  Finally scale a fixed smooth positive weight, equal to
$r_y^{-4}$ far out, so that $\rho\leq M_0/4$ everywhere.
This is possible by compact positivity and
$r_y^{-4}/r_y^{-3-\delta_*}=r_y^{-1+\delta_*}\to0$.
This proves the asserted uniform inequality for all $h$ and the
parameter order.
\end{proof}

Since $M_0\leq8\pi(\mu-|J|_g)$,
Proposition~\ref{ch:prep:prepared} satisfies the exact neutral
smallness hypothesis of Proposition~\ref{n3:domains:prepared-domain}.  The proof above
establishes the stronger charged choice directly.  At the assigned
face heights it also gives the exact boundary identities
\begin{equation}\label{ch:prep:face-values}
 h+C=c_b=P_B+H\quad(B_b),
 \qquad h+C=-c_w=P_B-H\quad(B_w),
\end{equation}
which will determine the signs in the neutral elliptic boundary
conditions.
